\documentclass[11pt]{article}
\usepackage{amsmath,amssymb,amsthm}
\usepackage[a4paper,margin=1in]{geometry}
\usepackage{enumitem}

\title{Statistics and Machine Learning Interview Notes}
\author{}
\date{}

\begin{document}
\maketitle

\section{Questions and Answers}

\subsection*{Q1. What is the Central Limit Theorem (CLT) and why is it useful?}

\textbf{Answer}

The Central Limit Theorem states that if
\[
X_1,X_2,\dots,X_n
\]
are i.i.d. random variables with mean $\mu$ and variance $\sigma^2$, then the sample mean
\[
\bar X=\frac{1}{n}\sum_{i=1}^{n}X_i
\]
satisfies
\[
\frac{\bar X-\mu}{\sigma/\sqrt n}
\overset{d}{\longrightarrow}
N(0,1)
\]
as $n\to\infty$. Equivalently,
\[
\bar X \approx N\!\left(\mu,\frac{\sigma^2}{n}\right).
\]

\textbf{Intuition}

Even if the original data distribution is highly skewed or non-normal, the distribution of sample means becomes approximately normal as the sample size increases.

\textbf{Example}

Consider the population
\[
\{1,1,1,1,100\}.
\]
This population is clearly not normal. However, if we repeatedly draw large samples and compute their means, the distribution of those means becomes approximately normal.

\textbf{Why is this useful?}

The CLT allows us to replace a complicated unknown distribution with a normal distribution when working with averages. Once something is approximately normal, we can compute probabilities, confidence intervals, and p-values.

\textbf{Example}

Suppose a sample of 100 people has average height
\[
\bar X = 171.2 \text{ cm}.
\]
Without the CLT, we would need the exact distribution of $\bar X$ to quantify uncertainty. The CLT tells us that
\[
\bar X \approx N\!\left(\mu,\frac{\sigma^2}{100}\right),
\]
allowing confidence intervals and hypothesis tests to be constructed.

\textbf{Connection to Hypothesis Testing}

For the z-test,
\[
Z=\frac{\bar X-\mu_0}{\sigma/\sqrt n}.
\]
Under the null hypothesis
\[
H_0:\mu=\mu_0,
\]
\[
Z \approx N(0,1),
\]
because of the CLT. This makes p-value computation possible.

\textbf{Machine Learning Example}

Let
\[
d_i= \text{PLTE contribution} - \text{SPE contribution}.
\]
Even if the individual $d_i$ values are highly non-normal, the average improvement
\[
\bar d=\frac{1}{n}\sum_{i=1}^{n} d_i
\]
is approximately normal for large $n$. Therefore, we can construct confidence intervals such as
\[
\bar d \pm 1.96\frac{s}{\sqrt n}.
\]

\textbf{Key Practical Benefit}

The CLT is the foundation of:
\begin{itemize}
    \item Confidence intervals
    \item z-tests
    \item t-tests
    \item p-values
    \item Sample size calculations
\end{itemize}

\textbf{In one sentence}

The Central Limit Theorem allows averages of large samples to be treated as normally distributed, enabling uncertainty quantification and hypothesis testing even when the original data is not normal.

\vspace{5mm}
\hrule
\vspace{5mm}

\subsection*{Q2. What is the advantage of a z-test over a t-test?}

\textbf{Answer}

The z-test and t-test are very similar. The main difference is that the z-test assumes the population standard deviation \(\sigma\) is known, whereas the t-test uses the sample standard deviation \(s\) because \(\sigma\) is usually unknown.

The test statistics are
\[
Z=\frac{\bar X-\mu_0}{\sigma/\sqrt n}
\]
and
\[
T=\frac{\bar X-\mu_0}{s/\sqrt n}.
\]

\textbf{Advantage of the z-test}

When \(\sigma\) is known, the denominator contains no estimation error. As a result:

\begin{itemize}
    \item The test statistic is more precise.
    \item The normal distribution has lighter tails than the t-distribution.
    \item The rejection threshold is slightly smaller.
    \item The test has slightly higher statistical power.
\end{itemize}

For example, for a two-sided \(5\%\) significance test,
\[
|Z| > 1.96
\]
whereas for a t-test with \(10\) degrees of freedom,
\[
|T| > 2.228.
\]

Because the t-test must account for the uncertainty in estimating \(\sigma\), stronger evidence is required to reject the null hypothesis.

\textbf{Why is the t-test used more often?}

In practice, the true population standard deviation \(\sigma\) is rarely known. Therefore, we typically estimate it using the sample standard deviation \(s\), which leads naturally to the t-test.

\textbf{Large Sample Behaviour}

As the sample size increases,
\[
s \approx \sigma
\]
and
\[
t_{n-1} \rightarrow N(0,1).
\]
Thus, for large samples, the t-test and z-test become nearly identical.

\textbf{Typical Usage}

\begin{itemize}
    \item \textbf{z-test:} Population standard deviation known.
    \item \textbf{t-test:} Population standard deviation unknown.
\end{itemize}

\textbf{In one sentence}

A z-test is slightly more powerful than a t-test because it uses the true population standard deviation, but since this quantity is rarely known, the t-test is the method used in most real-world applications.

\vspace{5mm}
\hrule
\vspace{5mm}

\subsection*{Q3. What is a confidence interval and why is it \(\bar X \pm 1.96\,\frac{s}{\sqrt n}\) for a normal distribution?}

\textbf{Answer}

A confidence interval (CI) is a range of values that is likely to contain the true population parameter.

For example, if we estimate a population mean \(\mu\) using the sample mean \(\bar X\), a \(95\%\) confidence interval provides a range of plausible values for \(\mu\).

The idea comes from the Central Limit Theorem, which states that for sufficiently large \(n\),
\[
\bar X \approx N\!\left(\mu,\frac{\sigma^2}{n}\right).
\]

Therefore,
\[
\frac{\bar X-\mu}{\sigma/\sqrt n}
\approx N(0,1).
\]

For a standard normal random variable \(Z\),
\[
P(-1.96 \le Z \le 1.96)=0.95.
\]

This means that \(95\%\) of the probability mass lies within \(1.96\) standard deviations of the mean.

Substituting
\[
Z=\frac{\bar X-\mu}{\sigma/\sqrt n},
\]
we obtain
\[
P\!\left(
-1.96
\le
\frac{\bar X-\mu}{\sigma/\sqrt n}
\le
1.96
\right)
=
0.95.
\]

Multiplying by \(\sigma/\sqrt n\),
\[
P\!\left(
\bar X-1.96\frac{\sigma}{\sqrt n}
\le
\mu
\le
\bar X+1.96\frac{\sigma}{\sqrt n}
\right)
=
0.95.
\]

Thus the \(95\%\) confidence interval is
\[
\boxed{
\bar X \pm 1.96\frac{\sigma}{\sqrt n}
}
\]
when the population standard deviation \(\sigma\) is known.

Since \(\sigma\) is usually unknown, we replace it with the sample standard deviation \(s\), giving
\[
\boxed{
\bar X \pm 1.96\frac{s}{\sqrt n}
}
\]
for large samples.

For smaller samples, the value \(1.96\) is replaced by the appropriate critical value from the t-distribution:
\[
\boxed{
\bar X \pm t_{\alpha/2,n-1}\frac{s}{\sqrt n}
}.
\]

\textbf{Interpretation}

A common misunderstanding is that there is a \(95\%\) probability that the true mean lies inside a computed confidence interval.

This is not correct.

The true mean \(\mu\) is fixed; the interval is random.

The correct interpretation is:

\begin{quote}
If we repeatedly draw samples and construct a \(95\%\) confidence interval from each sample, then approximately \(95\%\) of those intervals will contain the true population mean.
\end{quote}

\textbf{Example}

Suppose
\[
\bar X = 171.2,\qquad s=10,\qquad n=100.
\]

Then
\[
\text{SE}
=
\frac{s}{\sqrt n}
=
\frac{10}{10}
=
1.
\]

The \(95\%\) confidence interval is
\[
171.2 \pm 1.96(1)
=
171.2 \pm 1.96,
\]
or
\[
[169.24,\;173.16].
\]

\textbf{In one sentence}

A confidence interval is an estimate plus or minus a margin of error, and the factor \(1.96\) appears because \(95\%\) of a standard normal distribution lies between \(-1.96\) and \(+1.96\).

\vspace{5mm}
\hrule
\vspace{5mm}

\subsection*{Q4. What is the coefficient of variation (CV), and is it a statistical test?}

\textbf{Answer}

The coefficient of variation is a descriptive statistic that expresses the spread of a variable as a fraction of its average level:
\[
\boxed{
\mathrm{CV}=\frac{\sigma}{\mu}
}
\qquad\text{(sample version: } \mathrm{CV}=\frac{s}{\bar X}\text{)}.
\]
Because the units of \(\sigma\) and \(\mu\) cancel, the CV is unitless and can be used to compare variability across quantities measured on different scales.

\textbf{Example}

Suppose we train a classifier with 100 different random seeds and record the AUCPR each time. If
\[
\bar X = 0.60,\qquad s = 0.02,
\]
then
\[
\mathrm{CV}=\frac{0.02}{0.60}\approx 0.033,
\]
i.e.\ a typical seed's result deviates from the average by about \(3.3\%\).

\textbf{Is the CV a test?}

No. The CV is a descriptive statistic, not a hypothesis test:
\begin{itemize}
    \item There is no null hypothesis.
    \item There is no test statistic with a known sampling distribution under \(H_0\).
    \item There is no p-value.
\end{itemize}

\textbf{Limitations}

\begin{itemize}
    \item \textbf{No principled threshold.} Whether a CV of \(3\%\) or \(10\%\) is ``too large'' depends on the application, and the CV itself cannot tell us.
    \item \textbf{Unstable near zero.} If \(\mu\) is small (e.g.\ AUCPR \(\approx 0.05\)), even a tiny \(\sigma\) produces a large CV.
    \item \textbf{Measures spread, not bias.} A small CV only says that results are stable; it says nothing about whether the results are close to the \emph{correct} value.
    \item \textbf{Only meaningful for ratio-scale, positive data} (a true zero and no negative values).
\end{itemize}

\textbf{Do not confuse with cross-validation}

``CV'' is also the standard abbreviation for \emph{cross-validation}, which is an entirely different idea: the data is split into \(k\) folds, the model is trained on \(k-1\) folds and evaluated on the held-out fold, and the procedure is rotated through all folds. Cross-validation is a method for estimating model performance; the coefficient of variation is a summary of dispersion.

\textbf{In one sentence}

The coefficient of variation is the standard deviation divided by the mean, a unitless descriptive measure of relative spread (not a hypothesis test), and it has no built-in threshold for deciding when variability is ``too large''.

\vspace{5mm}
\hrule
\vspace{5mm}

\subsection*{Q5. How can we statistically test whether random undersampling of the majority class is ``good enough'' to capture the patterns in the full data?}

\textbf{Setting}

We have a very large, imbalanced dataset. The hypothesis is that a random undersample of the majority class is representative enough that a classifier trained on it performs essentially as well as one that has seen all of the majority data. We measure performance using AUCPR.

\textbf{A common first attempt: seed variability}

Draw the majority sample with \(100\) different seeds, train identical classifiers, obtain \(100\) AUCPR values, and judge the result by whether \(s/\bar X\) (the CV from Q4) is small.

This measures \emph{stability}: how much the result depends on which sample was drawn. It does \emph{not} measure \emph{representativeness}. If undersampling systematically discards a rare but important part of the majority class (for example hard negatives close to the decision boundary), then every one of the \(100\) samples loses it, and we obtain a tight distribution around a consistently mediocre AUCPR. Low variance can coexist with a systematic bias. In addition, there is no principled cutoff for ``too large''.

\textbf{Why a plain p-value test is also not the right tool}

Testing
\[
H_0:\ \mathrm{AUCPR}_{\text{full}}=\mathrm{AUCPR}_{\text{undersampled}}
\]
has several problems:
\begin{itemize}
    \item With enough seeds or enough test data, even a trivially small difference becomes ``significant'', and the result can be manufactured simply by running more seeds.
    \item Failing to reject \(H_0\) is not evidence of equivalence. ``Good enough'' is a claim of \emph{no meaningful loss}, which we need to show positively.
    \item The \(100\) runs share the same test set and the same minority data, so they are not independent draws from ``the population''. Their spread captures only majority-sampling noise, not test-set noise.
\end{itemize}

\textbf{Why a learning curve over the majority size is confounded}

Plotting AUCPR against the number of majority samples kept (1:1, 5:1, 20:1, all) changes two things at once: the amount of majority information (which should help) and the class-imbalance ratio (which hurts the learner, and adds noisy samples). A decline at larger ratios therefore cannot be attributed to the sample being unrepresentative or to the model simply struggling with imbalance. The curve cannot be interpreted as evidence either way.

\textbf{Recommended procedure}

\begin{enumerate}
    \item \textbf{Choose a reference that holds the imbalance ratio fixed.} Compare a single \(1{:}1\) undersample against a balanced-bagging / EasyEnsemble model, which averages many \(1{:}1\) undersamples drawn without replacement until the whole majority class is covered. Every base learner faces the same imbalance, but the ensemble has seen all of the majority data.
    \item \textbf{Fix the test set at natural prevalence.} AUCPR is highly sensitive to class prevalence, so the test set must not be resampled to a different imbalance.
    \item \textbf{Use a paired comparison.} Evaluate both approaches on identical folds and define the per-fold difference
    \[
    \Delta_i=\mathrm{AUCPR}^{\text{ref}}_i-\mathrm{AUCPR}^{\text{undersampled}}_i .
    \]
    \item \textbf{Test for equivalence or non-inferiority (TOST), not for difference.} Pre-specify a margin \(\delta\) that represents the largest acceptable loss (e.g.\ \(0.01\) AUCPR, chosen by what the business or fraud team cares about). With \(\Delta=\mathbb E[\Delta_i]\), the non-inferiority hypotheses are
    \[
    H_0:\ \Delta \ge \delta
    \qquad\text{vs.}\qquad
    H_1:\ \Delta < \delta .
    \]
    Rejecting \(H_0\) is positive evidence that undersampling is good enough. For a two-sided equivalence claim (TOST), both \(H_0:\Delta\ge\delta\) and \(H_0:\Delta\le-\delta\) must be rejected at level \(\alpha\). Equivalently, the \((1-2\alpha)\) confidence interval for \(\Delta\) (the \(90\%\) CI when \(\alpha=0.05\)) must lie entirely inside \((-\delta,\delta)\).
    \item \textbf{Obtain honest uncertainty.} Use repeated stratified \(k\)-fold cross-validation, or a bootstrap of the test set, to get a confidence interval for \(\Delta\). Because cross-validation estimates are correlated (training sets overlap), use the Nadeau--Bengio corrected resampled t-test, which inflates the variance of the mean difference:
    \[
    \widehat{\mathrm{Var}}(\bar\Delta)
    =
    \left(\frac{1}{J}+\frac{n_{\text{test}}}{n_{\text{train}}}\right)s_\Delta^{2},
    \]
    where \(J\) is the total number of resampling rounds and \(s_\Delta^{2}\) is the sample variance of the differences. The \(100\) seeds supply the majority-sampling component of the uncertainty, and the resampling supplies the test-set component.
    \item \textbf{Decision rule.} If the confidence interval for \(\Delta\) lies inside the pre-specified margin, conclude that random undersampling is good enough. Otherwise, it is not demonstrated to be.
\end{enumerate}

\textbf{Role of the seed-variance plot}

The distribution of AUCPR across the \(100\) seeds (and its CV) remains useful, but only as a \emph{secondary} stability diagnostic. The \emph{primary} evidence is the paired difference against the reference model, with a confidence interval and a pre-specified margin \(\delta\).

\textbf{Conceptual note}

A uniform random sample of the majority class is representative of the majority \emph{distribution} by construction. The real risk is \emph{coverage}: rare majority points near the decision boundary (hard negatives) matter most for precision at high recall, and a small sample may drop most of them. The ensemble-versus-single-sample comparison tests exactly this, because the ensemble recovers those points across its bags.

\textbf{In one sentence}

To test whether random undersampling is good enough, compare a single undersample against an ensemble covering the full majority class at the same imbalance ratio, using a paired, fixed-test-set comparison with a confidence interval, and conclude ``good enough'' only if the entire interval for the AUCPR difference lies within a pre-specified equivalence margin \(\delta\), rather than judging by the CV of seed results or a plain p-value.

\vspace{5mm}
\hrule
\vspace{5mm}

\subsection*{Q6. What is the variance of the sample mean \(\bar X=\frac{1}{n}\sum_{i=1}^{n}X_i\)? Derive it from the definition of variance.}

\textbf{Answer}

Let \(X_1,\dots,X_n\) be i.i.d. with \(E[X_i]=\mu\) and \(\mathrm{Var}(X_i)=\sigma^2\). We start from the definition
\[
\mathrm{Var}(Y)=E\big[(Y-E[Y])^2\big]
\]
and apply it to \(Y=\bar X\).

\textbf{Step 1: the mean of \(\bar X\).}
\[
E[\bar X]=\frac{1}{n}\sum_{i=1}^{n}E[X_i]=\frac{1}{n}\cdot n\mu=\mu.
\]

\textbf{Step 2: the deviation from the mean.}
\[
\bar X-E[\bar X]
=\frac{1}{n}\sum_{i=1}^{n}X_i-\mu
=\frac{1}{n}\sum_{i=1}^{n}(X_i-\mu),
\]
where we used \(\mu=\frac{1}{n}\sum_{i=1}^{n}\mu\) to bring \(\mu\) inside the sum.

\textbf{Step 3: substitute into the definition and square.}
\[
\mathrm{Var}(\bar X)
=E\!\left[\left(\frac{1}{n}\sum_{i=1}^{n}(X_i-\mu)\right)^{2}\right]
=\frac{1}{n^{2}}\,E\!\left[\sum_{i=1}^{n}\sum_{j=1}^{n}(X_i-\mu)(X_j-\mu)\right].
\]
Squaring a sum produces a double sum over all ordered pairs \((i,j)\).

\textbf{Step 4: move the expectation inside (linearity).}
\[
\mathrm{Var}(\bar X)
=\frac{1}{n^{2}}\sum_{i=1}^{n}\sum_{j=1}^{n}E\big[(X_i-\mu)(X_j-\mu)\big]
=\frac{1}{n^{2}}\sum_{i=1}^{n}\sum_{j=1}^{n}\mathrm{Cov}(X_i,X_j).
\]

\textbf{Step 5: split into diagonal and off-diagonal terms.}
\[
\mathrm{Var}(\bar X)
=\frac{1}{n^{2}}\left[
\sum_{i=1}^{n}\underbrace{E\big[(X_i-\mu)^2\big]}_{=\,\mathrm{Var}(X_i)\,=\,\sigma^{2}}
+\sum_{i\neq j}E\big[(X_i-\mu)(X_j-\mu)\big]
\right].
\]

\textbf{Step 6: use independence on the off-diagonal terms.}

For \(i\neq j\), independence makes the expectation of the product factor:
\[
E\big[(X_i-\mu)(X_j-\mu)\big]
=E[X_i-\mu]\;E[X_j-\mu]
=0\cdot 0
=0.
\]
So every cross term vanishes. The justification for the factorisation is given in detail below.

\textbf{Step 7: finish.}
\[
\mathrm{Var}(\bar X)
=\frac{1}{n^{2}}\sum_{i=1}^{n}\sigma^{2}
=\frac{n\sigma^{2}}{n^{2}}
=\boxed{\frac{\sigma^{2}}{n}}.
\]

\textbf{Why does independence make the expectation of the product factor?}

The key fact is the following. If \(U\) and \(V\) are independent random variables (with finite expectations), then
\[
E[UV]=E[U]\,E[V].
\]

\emph{Proof (continuous case).} Independence means the joint density is the product of the marginal densities:
\[
f_{U,V}(u,v)=f_U(u)\,f_V(v).
\]
Then, by the definition of expectation for a function of two random variables,
\[
E[UV]
=\iint uv\,f_{U,V}(u,v)\,du\,dv
=\iint uv\,f_U(u)\,f_V(v)\,du\,dv.
\]
Because the integrand is now a product of a function of \(u\) alone and a function of \(v\) alone, the double integral separates into a product of two single integrals:
\[
E[UV]
=\left(\int u\,f_U(u)\,du\right)\left(\int v\,f_V(v)\,dv\right)
=E[U]\,E[V].
\]

\emph{Discrete case.} The same argument holds with sums, using \(P(U=u,V=v)=P(U=u)P(V=v)\):
\[
E[UV]=\sum_{u}\sum_{v}uv\,P(U=u)P(V=v)
=\Big(\sum_u uP(U=u)\Big)\Big(\sum_v vP(V=v)\Big)=E[U]E[V].
\]

\emph{Applying it to our problem.} Take \(U=X_i-\mu\) and \(V=X_j-\mu\) with \(i\neq j\). Since \(U\) is a function of \(X_i\) only and \(V\) is a function of \(X_j\) only, and \(X_i\) and \(X_j\) are independent, \(U\) and \(V\) are also independent. Therefore
\[
E[(X_i-\mu)(X_j-\mu)]=E[X_i-\mu]\,E[X_j-\mu],
\]
and each factor is
\[
E[X_i-\mu]=E[X_i]-\mu=\mu-\mu=0.
\]
So the product is \(0\). Equivalently, \(\mathrm{Cov}(X_i,X_j)=0\) for \(i\neq j\).

\emph{What exactly is needed?} Only the weaker condition that the \(X_i\) are \emph{uncorrelated} (\(\mathrm{Cov}(X_i,X_j)=0\)) is required for the result \(\sigma^2/n\). Independence is a sufficient way to guarantee it. The converse does not hold: uncorrelated variables need not be independent.

\textbf{What changes if the observations are correlated?}

Suppose \(\mathrm{Cov}(X_i,X_j)=\rho\sigma^{2}\) for all \(i\neq j\). There are \(n(n-1)\) ordered off-diagonal pairs, so Step 5 gives
\[
\mathrm{Var}(\bar X)
=\frac{1}{n^{2}}\Big[n\sigma^{2}+n(n-1)\rho\sigma^{2}\Big]
=\frac{\sigma^{2}}{n}\Big[1+(n-1)\rho\Big].
\]
For \(\rho>0\) this is larger than \(\sigma^{2}/n\), and as \(n\to\infty\) it tends to \(\rho\sigma^{2}\) rather than to \(0\). Averaging correlated observations cannot remove their shared component of noise. This is why the Nadeau--Bengio correction in Q5 inflates the variance of the mean difference across overlapping cross-validation folds.

\textbf{Connection to earlier questions}

Taking the square root gives the standard error
\[
\mathrm{SE}(\bar X)=\frac{\sigma}{\sqrt n},
\]
which is the denominator in the CLT statement (Q1), the z- and t-statistics (Q2), and the confidence interval \(\bar X\pm1.96\,\sigma/\sqrt n\) (Q3). It also shows that uncertainty shrinks like \(1/\sqrt n\): quadrupling the sample size only halves the standard error.

\textbf{In one sentence}

For independent (or merely uncorrelated) observations with common variance \(\sigma^{2}\), the variance of the sample mean is \(\sigma^{2}/n\), because the cross-covariance terms vanish when expectations of products of independent variables factor into products of expectations, each of which is zero.

\vspace{5mm}
\hrule
\vspace{5mm}

\subsection*{Q7. What is a Z-test, and why does \(|Z|>1.96\) correspond to 5\% significance? Show mathematically.}

\textbf{Answer}

\textbf{What a Z-test is}

A Z-test checks whether a sample mean is far enough from a hypothesised value \(\mu_0\) that chance alone is an unlikely explanation. It assumes the population standard deviation \(\sigma\) is known, and that either the data are normal or \(n\) is large (so the CLT applies).
\[
H_0:\mu=\mu_0
\qquad\text{vs}\qquad
H_1:\mu\neq\mu_0,
\]
\[
Z=\frac{\bar X-\mu_0}{\sigma/\sqrt n}.
\]

Under \(H_0\), we have \(E[\bar X]=\mu_0\) and \(\mathrm{Var}(\bar X)=\sigma^2/n\) (from Q6). Standardising a variable (subtracting its mean and dividing by its standard deviation) gives mean \(0\) and variance \(1\):
\[
E[Z]=0,\qquad \mathrm{Var}(Z)=1,\qquad Z\sim N(0,1)
\]
(exactly if the data are normal, approximately by the CLT otherwise).

\textbf{Where 5\% and 1.96 come from}

\emph{Step 1: define the error rate we accept.} The significance level \(\alpha\) is the probability of rejecting \(H_0\) when it is true (Type I error). We choose a rejection region \(|Z|>c\) such that
\[
P(|Z|>c\mid H_0)=\alpha=0.05.
\]

\emph{Step 2: write this with the standard normal density.}
\[
\varphi(z)=\frac{1}{\sqrt{2\pi}}e^{-z^{2}/2},
\qquad
\Phi(c)=\int_{-\infty}^{c}\varphi(z)\,dz.
\]

\emph{Step 3: use symmetry.} Since \(\varphi\) is symmetric about \(0\), the two tails have equal probability:
\[
P(|Z|>c)=P(Z>c)+P(Z<-c)=2\,P(Z>c)=2\big(1-\Phi(c)\big).
\]

\emph{Step 4: solve for \(c\).}
\[
2\big(1-\Phi(c)\big)=0.05
\;\Longrightarrow\;
\Phi(c)=0.975
\;\Longrightarrow\;
c=\Phi^{-1}(0.975).
\]

\emph{Step 5: evaluate.} The integral \(\Phi(c)\) has no closed form, so \(c\) is obtained from numerical integration or tables:
\[
\Phi(1.96)=\int_{-\infty}^{1.96}\frac{1}{\sqrt{2\pi}}e^{-z^{2}/2}\,dz\approx 0.9750,
\]
so
\[
\boxed{c=1.96}
\qquad\text{and}\qquad
P(|Z|>1.96)=2(1-0.9750)=0.05.
\]

\textbf{Important clarification}

The event \(Z>1.96\) on its own has probability \(2.5\%\), not \(5\%\). The \(5\%\) figure applies to the \emph{two-sided} event \(|Z|>1.96\), i.e.\ \(Z>1.96\) \emph{or} \(Z<-1.96\), with each tail contributing \(2.5\%\).

\begin{center}
\begin{tabular}{lcc}
\hline
Test & Rejection region & Critical value \\
\hline
Two-sided, \(\alpha=0.05\) & \(|Z|>c\) & \(c=1.96\) \\
One-sided, \(\alpha=0.05\) & \(Z>c\) & \(c=1.645\) \\
\hline
\end{tabular}
\end{center}

For the one-sided case, \(1-\Phi(c)=0.05\Rightarrow\Phi(c)=0.95\Rightarrow c=1.645\).

\textbf{The p-value version}

Given an observed value \(z_{\text{obs}}\),
\[
p=P\big(|Z|\ge|z_{\text{obs}}|\mid H_0\big)=2\big(1-\Phi(|z_{\text{obs}}|)\big).
\]
Rejecting when \(|z_{\text{obs}}|>1.96\) is equivalent to rejecting when \(p<0.05\), because \(p\) is a decreasing function of \(|z_{\text{obs}}|\) and equals \(0.05\) exactly at \(|z_{\text{obs}}|=1.96\).

\textbf{Example}

Let \(\bar X=171.2\), \(\mu_0=170\), \(\sigma=10\), \(n=100\). Then
\[
z=\frac{171.2-170}{10/\sqrt{100}}=1.2,
\qquad
p=2\big(1-\Phi(1.2)\big)=2(1-0.8849)=0.230.
\]
Since \(|1.2|<1.96\) (equivalently \(p>0.05\)), we fail to reject \(H_0\).

\textbf{Link to the confidence interval (Q3)}

The acceptance region of the test is \(|Z|\le1.96\):
\[
-1.96\le\frac{\bar X-\mu_0}{\sigma/\sqrt n}\le1.96
\iff
\mu_0\in\left[\bar X-1.96\frac{\sigma}{\sqrt n},\ \bar X+1.96\frac{\sigma}{\sqrt n}\right].
\]
So the \(95\%\) confidence interval is exactly the set of values \(\mu_0\) that the two-sided \(5\%\) Z-test would not reject, which is why the same constant \(1.96\) appears in both.

\textbf{In one sentence}

A Z-test standardises the sample mean under \(H_0\) so that it follows \(N(0,1)\), and the two-sided critical value \(1.96\) is the point satisfying \(P(|Z|>1.96)=2(1-\Phi(1.96))=0.05\), because \(2.5\%\) of the standard normal probability lies in each tail beyond \(\pm1.96\).

\end{document}